\clearpage
\section{Atom equality by resolved store identity}
\label{sec:copying-atom-equality}
Both Figure~19(d) rules (PDF p.~38) now have concrete finite-graph
implementations. Two instances attached to one atom contraction tree and
store compare true. Instances attached to two separate trees and stores
compare false. Both instances are weakened after the comparison, while all
trees, stores and store-content interfaces remain. The trees and the
numbers of other references are arbitrary.

The decisive semantic check is identity of the resolved store, not equality
of stored contents. Successful lookups determine a unique boolean even in
cyclic hosts, and inserting a fragment into a surrounding graph preserves
that boolean. Each of the two local rules has a checked seven-step
search/compute/return execution. Renaming, forward padding and recursive
label safety are proved. These results complete the local constructions
for all four Figure~19 store operations.

They do not yet prove determinism of the combined copying UAM, arbitrary
focused-context matching, or the robustness of the templates. The lookup
outcome uniqueness theorem below is narrower than uniqueness of an entire
graph rewrite. Proposition~12.2 and its Figure~16 translation remain
unchanged, and the full proposition is still unfinished.

\Needspace{12\baselineskip}
\subsection{The common-store frame}

\termdef{copying-equality-same}
The common-store fragment has one external term input, $h$ other atom inputs,
and one term output for the contents of the shared store. Its tree slot has
$h+2$ atom inputs: the first two come from the two instances; the rest are
the other references. The redex edge tags are comparison, left instance,
right instance, tree and store. The reduct replaces the comparison by a
true constant and adds a separate weakening for each instance.

\lean{Copying/AtomEqualitySameFrames}{8}{8}{AtomEquality.Same.inputs}

\lean{Copying/AtomEqualitySameFrames}{9}{9}{AtomEquality.Same.treeHole}

\lean{Copying/AtomEqualitySameFrames}{10}{12}{AtomEquality.Same.BeforeEdge}

\lean{Copying/AtomEqualitySameFrames}{13}{15}{AtomEquality.Same.AfterEdge}

The external input index is unchanged by the boundary embedding. The
other-reference list consists of successor positions in increasing order.
Injectivity of successor and the input constructor proves that this list
has no duplicates; finite-range membership characterises its elements.

\lean{Copying/AtomEqualitySameFrames}{17}{18}{AtomEquality.Same.inputPort}

\lean{Copying/AtomEqualitySameFrames}{19}{20}{AtomEquality.Same.otherInputs}

\lean{Copying/AtomEqualitySameFrames}{21}{22}{AtomEquality.Same.otherInputs_nodup}

\lean{Copying/AtomEqualitySameFrames}{23}{24}{AtomEquality.Same.mem_otherInputs}

Five internal vertices carry the wiring: two term argument inputs, two
atom outputs of the instances, and the atom connection between tree root
and store. The comparison's ordered outputs feed the two term vertices.
The tree's first two inputs receive the two atom vertices; its root feeds
the store; the store's term output is the external output. Label, source
and target functions spell out this graph. The constructor proves finiteness,
complete vertex coverage, unique incidence, duplicate-free lists and
signature typing by finite case analysis and the indexed other-input list.

\lean{Copying/AtomEqualitySameFrames}{26}{31}{AtomEquality.Same.beforeLabel}

\lean{Copying/AtomEqualitySameFrames}{33}{38}{AtomEquality.Same.beforeSource}

\lean{Copying/AtomEqualitySameFrames}{40}{45}{AtomEquality.Same.beforeTarget}

\lean{Copying/AtomEqualitySameFrames}{47}{98}{AtomEquality.Same.beforeFrame}

On the reduct side, each consumed instance receives a term weakening at
its original term input. The true constant consumes the external term
input and has no output. The atom tree and store wiring are identical to
the redex wiring. Seven finite edge cases establish the same coverage,
uniqueness and typing obligations; no tree or store edge is discarded.

\lean{Copying/AtomEqualitySameFrames}{100}{107}{AtomEquality.Same.afterLabel}

\lean{Copying/AtomEqualitySameFrames}{109}{116}{AtomEquality.Same.afterSource}

\lean{Copying/AtomEqualitySameFrames}{118}{125}{AtomEquality.Same.afterTarget}

\lean{Copying/AtomEqualitySameFrames}{127}{178}{AtomEquality.Same.afterFrame}

\Needspace{12\baselineskip}
\subsection{The distinct-store frame}

\termdef{copying-equality-different}
The distinct-store fragment has one external term input, $n+m$ other atom
inputs, and two term outputs for the two stores' contents. Each tree slot
has one selected instance input and its own other references: $n$ on the
left, $m$ on the right. The two tree holes have different names. The edge
types explicitly contain two store occurrences, in addition to two trees
and two instances. Including these stores is essential: two separate
reference paths alone would not prove that their eventual stores differ.

\lean{Copying/AtomEqualityDifferentFrames}{8}{8}{AtomEquality.Different.inputs}

\lean{Copying/AtomEqualityDifferentFrames}{9}{9}{AtomEquality.Different.treeHole}

\lean{Copying/AtomEqualityDifferentFrames}{10}{12}{AtomEquality.Different.BeforeEdge}

\lean{Copying/AtomEqualityDifferentFrames}{13}{15}{AtomEquality.Different.AfterEdge}

The left other-reference block occupies positions $1$ through $n$; the
right block occupies positions $n+1$ through $n+m$. The finite lists use
these offsets explicitly. Their duplicate-freedom proofs reduce equality
of port indices to equality of the corresponding finite indices. Membership
lemmas expose the index witnesses, allowing the frame proof to show that
the two blocks cannot overlap.

\lean{Copying/AtomEqualityDifferentFrames}{17}{18}{AtomEquality.Different.inputPort}

\lean{Copying/AtomEqualityDifferentFrames}{19}{20}{AtomEquality.Different.leftInputs}

\lean{Copying/AtomEqualityDifferentFrames}{21}{22}{AtomEquality.Different.rightInputs}

\lean{Copying/AtomEqualityDifferentFrames}{23}{24}{AtomEquality.Different.leftInputs_nodup}

\lean{Copying/AtomEqualityDifferentFrames}{25}{26}{AtomEquality.Different.rightInputs_nodup}

\lean{Copying/AtomEqualityDifferentFrames}{27}{28}{AtomEquality.Different.mem_leftInputs}

\lean{Copying/AtomEqualityDifferentFrames}{29}{30}{AtomEquality.Different.mem_rightInputs}

The redex has two internal term vertices and four internal atom vertices.
The two instance outputs enter their respective tree inputs; each tree
root feeds its own store; the stores feed external term outputs zero and
one. The other-reference blocks attach to the corresponding trees. The
net proof checks coverage by separating external input zero, the left
block and the right block. Source uniqueness uses the disjoint index
intervals; target uniqueness separates the two store outputs. The remaining
finite cases establish typing and all incidence requirements.

\lean{Copying/AtomEqualityDifferentFrames}{32}{39}{AtomEquality.Different.beforeLabel}

\lean{Copying/AtomEqualityDifferentFrames}{41}{48}{AtomEquality.Different.beforeSource}

\lean{Copying/AtomEqualityDifferentFrames}{50}{57}{AtomEquality.Different.beforeTarget}

\lean{Copying/AtomEqualityDifferentFrames}{59}{128}{AtomEquality.Different.beforeFrame}

The reduct retains those six internal vertices, both trees and both
stores. It weakens each consumed instance and places a false constant at
the external term input. The nine-edge frame proves the same interface
and typing properties. The two output positions remain distinct, so the
store identities cannot be merged merely by choosing equal contents.

\lean{Copying/AtomEqualityDifferentFrames}{130}{139}{AtomEquality.Different.afterLabel}

\lean{Copying/AtomEqualityDifferentFrames}{141}{150}{AtomEquality.Different.afterSource}

\lean{Copying/AtomEqualityDifferentFrames}{152}{161}{AtomEquality.Different.afterTarget}

\lean{Copying/AtomEqualityDifferentFrames}{163}{232}{AtomEquality.Different.afterFrame}

\Needspace{12\baselineskip}
\subsection{Literal rule graphs and depth allowances}

For the common-store case, the filling functions put the supplied arbitrary
atom contraction tree in its slot and fill every other edge with its
original primitive corolla. Expansion followed by primitive relabelling
produces the two actual hypernets at every allowance $d$. The rule has one
constructor relating exactly those graphs. It requires an actual contraction
tree, without any assumed machine determinism or observational equivalence.

\lean{Copying/AtomEquality}{8}{16}{AtomEquality.Same.beforeFill}

\lean{Copying/AtomEquality}{17}{19}{AtomEquality.Same.before}

\lean{Copying/AtomEquality}{20}{30}{AtomEquality.Same.afterFill}

\lean{Copying/AtomEquality}{31}{33}{AtomEquality.Same.after}

\lean{Copying/AtomEquality}{34}{35}{AtomEquality.Same.Rule}

The designated left instance, right instance and store are explicit edge
injections on each side. They provide the coordinates used by all lookup
and execution proofs. Raising the allowance preserves each rule graph
literally because all of its labels are primitive and the graph wiring is
unchanged. Rule preservation follows by rewriting these equalities; induction
gives every finite padding amount.

\lean{Copying/AtomEquality}{36}{36}{AtomEquality.Same.beforeLeft}

\lean{Copying/AtomEquality}{37}{37}{AtomEquality.Same.beforeRight}

\lean{Copying/AtomEquality}{38}{38}{AtomEquality.Same.beforeStore}

\lean{Copying/AtomEquality}{40}{40}{AtomEquality.Same.afterLeft}

\lean{Copying/AtomEquality}{41}{41}{AtomEquality.Same.afterRight}

\lean{Copying/AtomEquality}{42}{42}{AtomEquality.Same.afterStore}

\lean{Copying/AtomEquality}{39}{39}{AtomEquality.Same.before_raise}

\lean{Copying/AtomEquality}{43}{43}{AtomEquality.Same.after_raise}

\lean{Copying/AtomEquality}{44}{47}{AtomEquality.Same.rule_raise}

\lean{Copying/AtomEquality}{48}{52}{AtomEquality.Same.rule_pad}

The distinct-store filling functions put the two arbitrary trees into
separate slots, retaining the other primitive corollas. The rule likewise
relates the literal redex and reduct. Separate edge injections identify
both instances and both stores before and after the rewrite. Its primitive
construction has the same literal raising equalities and finite-padding
preservation, again proved directly from the constructors.

\lean{Copying/AtomEquality}{55}{65}{AtomEquality.Different.beforeFill}

\lean{Copying/AtomEquality}{66}{68}{AtomEquality.Different.before}

\lean{Copying/AtomEquality}{69}{81}{AtomEquality.Different.afterFill}

\lean{Copying/AtomEquality}{82}{84}{AtomEquality.Different.after}

\lean{Copying/AtomEquality}{85}{86}{AtomEquality.Different.Rule}

\lean{Copying/AtomEquality}{87}{87}{AtomEquality.Different.beforeLeft}

\lean{Copying/AtomEquality}{88}{88}{AtomEquality.Different.beforeRight}

\lean{Copying/AtomEquality}{89}{89}{AtomEquality.Different.beforeStore}

\lean{Copying/AtomEquality}{90}{90}{AtomEquality.Different.beforeStoreRight}

\lean{Copying/AtomEquality}{92}{92}{AtomEquality.Different.afterLeft}

\lean{Copying/AtomEquality}{93}{93}{AtomEquality.Different.afterRight}

\lean{Copying/AtomEquality}{94}{94}{AtomEquality.Different.afterStore}

\lean{Copying/AtomEquality}{95}{95}{AtomEquality.Different.afterStoreRight}

\lean{Copying/AtomEquality}{91}{91}{AtomEquality.Different.before_raise}

\lean{Copying/AtomEquality}{96}{96}{AtomEquality.Different.after_raise}

\lean{Copying/AtomEquality}{97}{100}{AtomEquality.Different.rule_raise}

\lean{Copying/AtomEquality}{101}{105}{AtomEquality.Different.rule_pad}

\Needspace{12\baselineskip}
\subsection{Resolving both references before and after comparison}

The common path lemma concerns one contraction tree embedded into two
actual hypernets. Suppose its adjacency steps, selected input attachment
and root-to-store attachment are preserved in both embeddings, and the
two designated destinations have store labels. Tree reachability supplies
a path of length $q$ from the selected input edge to the root. Map that
same path into both graphs, prepend the instance hop and append the store
hop. Both resolutions therefore have exactly $q+2$ hops. The concrete
applications below discharge every embedding premise; the local rules do
not assume these routing properties.

\lean{Copying/AtomEqualityResolution}{9}{22}{AtomEquality.tree_resolution_pair}

In the common-store case, the left and right tree-edge injections use the
same tree slot, but the entry inputs are zero and one respectively. The
vertex maps are expansion insertion maps. Adjacent tree edges become atom
hops because their labels are weakening or contraction and their shared
vertex is preserved. The instance-target calculation proves entry, and
the tree-root output calculation proves exit to the shared store.

\lean{Copying/AtomEqualityResolution}{24}{24}{AtomEquality.Same.beforeLeftTree}

\lean{Copying/AtomEqualityResolution}{25}{26}{AtomEquality.Same.beforeLeftVertex}

\lean{Copying/AtomEqualityResolution}{27}{27}{AtomEquality.Same.afterLeftTree}

\lean{Copying/AtomEqualityResolution}{28}{29}{AtomEquality.Same.afterLeftVertex}

\lean{Copying/AtomEqualityResolution}{30}{68}{AtomEquality.Same.resolves_left}

\lean{Copying/AtomEqualityResolution}{69}{69}{AtomEquality.Same.beforeRightTree}

\lean{Copying/AtomEqualityResolution}{70}{71}{AtomEquality.Same.beforeRightVertex}

\lean{Copying/AtomEqualityResolution}{72}{72}{AtomEquality.Same.afterRightTree}

\lean{Copying/AtomEqualityResolution}{73}{74}{AtomEquality.Same.afterRightVertex}

\lean{Copying/AtomEqualityResolution}{75}{113}{AtomEquality.Same.resolves_right}

The resulting existence theorems preserve each reference's own exact
resolution count across the rewrite. They do not assert that the left
and right references have equal counts. For any proposed count, uniqueness
of successful store traces identifies it with the constructed count; this
proves the two exact-count equivalences.

\lean{Copying/AtomEqualityResolution}{114}{120}{AtomEquality.Same.left_resolution_count_iff}

\lean{Copying/AtomEqualityResolution}{121}{127}{AtomEquality.Same.right_resolution_count_iff}

In the distinct-store case, the left edge and vertex maps use the first
tree slot and the right maps use the second. Each selected entry is input
zero of its own tree. The two root outputs attach to their own designated
stores. The same adjacency, entry and exit calculations discharge the
common path lemma for each reference separately.

\lean{Copying/AtomEqualityResolution}{130}{130}{AtomEquality.Different.beforeLeftTree}

\lean{Copying/AtomEqualityResolution}{131}{132}{AtomEquality.Different.beforeLeftVertex}

\lean{Copying/AtomEqualityResolution}{133}{133}{AtomEquality.Different.afterLeftTree}

\lean{Copying/AtomEqualityResolution}{134}{135}{AtomEquality.Different.afterLeftVertex}

\lean{Copying/AtomEqualityResolution}{136}{174}{AtomEquality.Different.resolves_left}

\lean{Copying/AtomEqualityResolution}{175}{175}{AtomEquality.Different.beforeRightTree}

\lean{Copying/AtomEqualityResolution}{176}{177}{AtomEquality.Different.beforeRightVertex}

\lean{Copying/AtomEqualityResolution}{178}{178}{AtomEquality.Different.afterRightTree}

\lean{Copying/AtomEqualityResolution}{179}{180}{AtomEquality.Different.afterRightVertex}

\lean{Copying/AtomEqualityResolution}{181}{219}{AtomEquality.Different.resolves_right}

The resulting existence theorems preserve each reference's own exact
resolution count across the rewrite. They do not assert that the left
and right references have equal counts. For any proposed count, uniqueness
of successful store traces identifies it with the constructed count; this
proves the two exact-count equivalences.

\lean{Copying/AtomEqualityResolution}{220}{226}{AtomEquality.Different.left_resolution_count_iff}

\lean{Copying/AtomEqualityResolution}{227}{233}{AtomEquality.Different.right_resolution_count_iff}

\Needspace{12\baselineskip}
\subsection{Store identity determines a unique boolean}

\termdef{copying-equality-outcome}
For a graph $G$ and selected edges $a,b$, a justified comparison outcome
$r$ consists of successful store traces $a\leadsto s$ and $b\leadsto t$
with their individual lengths, together with $r=\mathsf{true}$ if and only
if $s=t$. The definition compares the destination edges themselves. It
never inspects their term outputs or stored contents.

Two such witnesses must have the same left store and the same right store,
by uniqueness of successful store resolution. Substituting these identities
leaves the same truth condition for both booleans, so Boolean case analysis
proves their equality. This result needs no global acyclicity and no
assumption that the full UAM is deterministic.

\lean{Copying/AtomEqualityIdentity}{9}{11}{AtomEquality.Outcome}

\lean{Copying/AtomEqualityIdentity}{14}{21}{AtomEquality.outcome_unique}

Inserting a matched fragment into any correctly typed hole preserves both
successful traces with their exact lengths. Store occurrences are mapped
by the inserted-edge injection, which preserves and reflects equality.
Consequently the comparison outcome is preserved. Applying outcome
uniqueness in the resulting host graph excludes any contrary successful
outcome obtained by a different routing witness. This proves locality of
the identity decision; it does not by itself prove contextual closure or
uniqueness of all graph rewrites.

\lean{Copying/AtomEqualityIdentity}{24}{30}{AtomEquality.outcome_insert}

\lean{Copying/AtomEqualityIdentity}{33}{37}{AtomEquality.inserted_outcome_unique}

In the common-store construction, the two resolution theorems have the
same destination edge, on both rule sides. They therefore justify true.
A false outcome is impossible by uniqueness of the comparison result.

\lean{Copying/AtomEqualityIdentity}{41}{45}{AtomEquality.Same.true_before}

\lean{Copying/AtomEqualityIdentity}{47}{51}{AtomEquality.Same.true_after}

\lean{Copying/AtomEqualityIdentity}{53}{55}{AtomEquality.Same.cannot_return_false}

In the distinct-store construction, equality of the designated stores
would force two distinct frame tags to be equal. This contradiction proves
store distinctness before and after the rule. Together with the two
successful resolutions it justifies false on both sides. A true outcome
is therefore impossible, regardless of the contents attached to the stores.

\lean{Copying/AtomEqualityIdentity}{59}{61}{AtomEquality.Different.stores_distinct_before}

\lean{Copying/AtomEqualityIdentity}{63}{65}{AtomEquality.Different.stores_distinct_after}

\lean{Copying/AtomEqualityIdentity}{67}{71}{AtomEquality.Different.false_before}

\lean{Copying/AtomEqualityIdentity}{73}{77}{AtomEquality.Different.false_after}

\lean{Copying/AtomEqualityIdentity}{79}{81}{AtomEquality.Different.cannot_return_true}

\Needspace{12\baselineskip}
\subsection{The search and equality subsystem}

\termdef{copying-equality-subsystem}
A state consists of its actual graph and token position. One common local
transition relation contains intrinsic search and both equality rewrites,
with the interfaces fixed by their respective fragments. Each rewrite
starts in rewrite mode at the external term input and resets search at
that input in the reduct. Finality is backtracking at an external input.
The search finality theorem excludes another search transition there;
both equality constructors require a different focus mode. Thus final
normality is proved for this subsystem machine.

\lean{Copying/AtomEqualityExecution}{8}{8}{AtomEquality.State}

\lean{Copying/AtomEqualityExecution}{9}{15}{AtomEquality.Step}

\lean{Copying/AtomEqualityExecution}{17}{27}{AtomEquality.final_normal}

\lean{Copying/AtomEqualityExecution}{29}{32}{AtomEquality.machine}

For the common-store case, the argument coordinates
are the two internal term vertices. Search enters the left instance, returns
from it, resumes at the right instance, and returns from it. Backtracking
through the comparison head then requests rewriting. The result-edge
selector identifies the literal true constant in the reduct, and its label
theorem checks that value directly. Searching that constant returns at the
external input in one step. Composing the five pre-compute search steps,
the equality rewrite and the final return gives exactly seven transitions.
These are local fragment executions; a complete execution of a translated
source equality expression in an arbitrary context is not claimed here.

\lean{Copying/AtomEqualityExecution}{34}{34}{AtomEquality.Same.leftArgument}

\lean{Copying/AtomEqualityExecution}{35}{35}{AtomEquality.Same.rightArgument}

\lean{Copying/AtomEqualityExecution}{36}{40}{AtomEquality.Same.search_left}

\lean{Copying/AtomEqualityExecution}{41}{44}{AtomEquality.Same.left_returns}

\lean{Copying/AtomEqualityExecution}{45}{48}{AtomEquality.Same.right_returns}

\lean{Copying/AtomEqualityExecution}{49}{51}{AtomEquality.Same.search_right}

\lean{Copying/AtomEqualityExecution}{53}{55}{AtomEquality.Same.ready}

\lean{Copying/AtomEqualityExecution}{57}{57}{AtomEquality.Same.resultEdge}

\lean{Copying/AtomEqualityExecution}{58}{59}{AtomEquality.Same.result_label}

\lean{Copying/AtomEqualityExecution}{61}{65}{AtomEquality.Same.result_returns}

\lean{Copying/AtomEqualityExecution}{67}{73}{AtomEquality.Same.comparison_run}

For the distinct-store case, the argument coordinates
are the two internal term vertices. Search enters the left instance, returns
from it, resumes at the right instance, and returns from it. Backtracking
through the comparison head then requests rewriting. The result-edge
selector identifies the literal false constant in the reduct, and its label
theorem checks that value directly. Searching that constant returns at the
external input in one step. Composing the five pre-compute search steps,
the equality rewrite and the final return gives exactly seven transitions.
These are local fragment executions; a complete execution of a translated
source equality expression in an arbitrary context is not claimed here.

\lean{Copying/AtomEqualityExecution}{76}{76}{AtomEquality.Different.leftArgument}

\lean{Copying/AtomEqualityExecution}{77}{77}{AtomEquality.Different.rightArgument}

\lean{Copying/AtomEqualityExecution}{78}{82}{AtomEquality.Different.search_left}

\lean{Copying/AtomEqualityExecution}{83}{86}{AtomEquality.Different.left_returns}

\lean{Copying/AtomEqualityExecution}{87}{90}{AtomEquality.Different.right_returns}

\lean{Copying/AtomEqualityExecution}{91}{93}{AtomEquality.Different.search_right}

\lean{Copying/AtomEqualityExecution}{95}{97}{AtomEquality.Different.ready}

\lean{Copying/AtomEqualityExecution}{99}{99}{AtomEquality.Different.resultEdge}

\lean{Copying/AtomEqualityExecution}{100}{101}{AtomEquality.Different.result_label}

\lean{Copying/AtomEqualityExecution}{103}{107}{AtomEquality.Different.result_returns}

\lean{Copying/AtomEqualityExecution}{109}{115}{AtomEquality.Different.comparison_run}

\Needspace{12\baselineskip}
\subsection{Representation and boundary preservation}

For the common-store construction, ordered tree-carrier
isomorphisms extend to both rule sides. Componentwise expansion uses the
supplied isomorphism in each tree slot and identity on every primitive
corolla. Primitive relabelling lifts that isomorphism to the required
allowance. The two weakening selectors then prove explicitly that each
consumed instance is fed by a term weakening at its original source.

\lean{Copying/AtomEqualityPreservation}{8}{16}{AtomEquality.Same.before_congr}

\lean{Copying/AtomEqualityPreservation}{17}{27}{AtomEquality.Same.after_congr}

\lean{Copying/AtomEqualityPreservation}{28}{28}{AtomEquality.Same.discardLeft}

\lean{Copying/AtomEqualityPreservation}{29}{31}{AtomEquality.Same.left_weakened}

\lean{Copying/AtomEqualityPreservation}{32}{32}{AtomEquality.Same.discardRight}

\lean{Copying/AtomEqualityPreservation}{33}{35}{AtomEquality.Same.right_weakened}

The single store target remains external output zero. Every other tree
input $i+2$ remains external input $i+1$ before and after comparison, as
shown by indexing the two-prefix source list. The stored content and all
other-reference attachments therefore retain their ordered boundary positions.

\lean{Copying/AtomEqualityPreservation}{36}{38}{AtomEquality.Same.store_boundary}

\lean{Copying/AtomEqualityPreservation}{39}{43}{AtomEquality.Same.other_references}

For the distinct-store construction, ordered tree-carrier
isomorphisms extend to both rule sides. Componentwise expansion uses the
supplied isomorphism in each tree slot and identity on every primitive
corolla. Primitive relabelling lifts that isomorphism to the required
allowance. The two weakening selectors then prove explicitly that each
consumed instance is fed by a term weakening at its original source.

\lean{Copying/AtomEqualityPreservation}{46}{56}{AtomEquality.Different.before_congr}

\lean{Copying/AtomEqualityPreservation}{57}{69}{AtomEquality.Different.after_congr}

\lean{Copying/AtomEqualityPreservation}{70}{70}{AtomEquality.Different.discardLeft}

\lean{Copying/AtomEqualityPreservation}{71}{73}{AtomEquality.Different.left_weakened}

\lean{Copying/AtomEqualityPreservation}{74}{74}{AtomEquality.Different.discardRight}

\lean{Copying/AtomEqualityPreservation}{75}{77}{AtomEquality.Different.right_weakened}

The two store targets remain external outputs zero and one respectively.
Other left-tree input $i+1$ remains input $i+1$; other right-tree input
$j+1$ remains input $n+j+1$. The finite-range source lists prove these
four before/after equalities without permuting reference blocks.

\lean{Copying/AtomEqualityPreservation}{78}{82}{AtomEquality.Different.store_boundaries}

\lean{Copying/AtomEqualityPreservation}{83}{87}{AtomEquality.Different.other_left_references}

\lean{Copying/AtomEqualityPreservation}{88}{92}{AtomEquality.Different.other_right_references}

\Needspace{12\baselineskip}
\subsection{Safety and stored-content regressions}

Recursive label safety for the common-store fragments follows
by classifying the finite frame edges. Tree labels satisfy the chosen
predicate whenever atom weakening and contraction do; the other edges
require the appropriate comparison/result operation, instances, stores
and term weakenings. Specialising the predicate to exclude focus labels,
or hole labels, proves both rule sides focus-free and hole-free at every
allowance. The temporary tree slots have been filled by actual graphs.

\lean{Copying/AtomEqualitySafety}{7}{18}{AtomEquality.Same.before_labels}

\lean{Copying/AtomEqualitySafety}{19}{32}{AtomEquality.Same.after_labels}

\lean{Copying/AtomEqualitySafety}{33}{38}{AtomEquality.Same.focusFree_sides}

\lean{Copying/AtomEqualitySafety}{39}{44}{AtomEquality.Same.holeFree_sides}

Recursive label safety for the distinct-store fragments follows
by classifying the finite frame edges. Tree labels satisfy the chosen
predicate whenever atom weakening and contraction do; the other edges
require the appropriate comparison/result operation, instances, stores
and term weakenings. Specialising the predicate to exclude focus labels,
or hole labels, proves both rule sides focus-free and hole-free at every
allowance. The temporary tree slots have been filled by actual graphs.

\lean{Copying/AtomEqualitySafety}{47}{61}{AtomEquality.Different.before_labels}

\lean{Copying/AtomEqualitySafety}{62}{78}{AtomEquality.Different.after_labels}

\lean{Copying/AtomEqualitySafety}{79}{84}{AtomEquality.Different.focusFree_sides}

\lean{Copying/AtomEqualitySafety}{85}{90}{AtomEquality.Different.holeFree_sides}

Serially attaching arbitrary contents preserves a justified outcome.
Both successful traces lift into the left serial component; equality of
store edges is preserved and reflected by its injection. The common-store
case thus compares true for arbitrary contents. A weakening-free binary
tree supplies a checked noncanonical seven-step true execution.

\lean{Copying/AtomEqualityExamples}{9}{17}{AtomEquality.outcome_serial_left}

\lean{Copying/AtomEqualityExamples}{21}{23}{AtomEquality.Same.arbitrary_contents}

\lean{Copying/AtomEqualityExamples}{25}{28}{AtomEquality.Same.noncanonical_run}

The distinct-store case compares false for arbitrary two-input contents.
No surrounding content graph can identify the designated stores in a
successful comparison: a hypothetical true outcome contradicts outcome
uniqueness. The theorem does not require the contents to terminate or the
surrounding graph to be acyclic. It is a routing theorem, not an assertion
that arbitrary callbacks have already been executed in the full language.

\lean{Copying/AtomEqualityExamples}{32}{35}{AtomEquality.Different.arbitrary_contents}

\lean{Copying/AtomEqualityExamples}{37}{41}{AtomEquality.Different.contents_cannot_identify_stores}

The final regression attaches two copies of the unchanged Figure~16
translation of the same natural numeral to the two stores. The heap and
literal-edge selectors retain their actual distributors and component
tags. Both literal labels denote the same numeral, and explicit incidence
equalities connect the respective store targets to those literal sources.
Nevertheless the reference comparison is false, by the arbitrary-contents
theorem. This checks that equality has not been replaced by comparison of
stored values.

\lean{Copying/AtomEqualityExamples}{44}{45}{AtomEquality.Different.identicalContents}

\lean{Copying/AtomEqualityExamples}{47}{49}{AtomEquality.Different.equalValueHeap}

\lean{Copying/AtomEqualityExamples}{51}{52}{AtomEquality.Different.equalValueLeft}

\lean{Copying/AtomEqualityExamples}{53}{54}{AtomEquality.Different.equalValueRight}

\lean{Copying/AtomEqualityExamples}{56}{58}{AtomEquality.Different.identical_literal_labels}

\lean{Copying/AtomEqualityExamples}{60}{62}{AtomEquality.Different.identical_literal_attachments}

\lean{Copying/AtomEqualityExamples}{65}{67}{AtomEquality.Different.identical_contents_compare_false}

The full build checks 335 jobs. All three axiom audits pass: 6 abstract,
477 graph and 652 copying results, using only standard Lean logical axioms.
Estimated completion of Proposition~12.2 is about 67\% by mathematical effort,
with broad uncertainty of 52--77\% (26 October 2026). All four store operations
now have local concrete rules. Arithmetic and application remain to be
constructed, followed by the combined-machine laws, template robustness,
semantic weakening and Figure~22 equivalence chain. Appendix~B.1 retains
those obligations and removes atom-equality construction from pending work.

